% Created 2026-11-23 Mon 23:23
% Intended LaTeX compiler: pdflatex
\documentclass[11pt]{article}
\usepackage[utf8]{inputenc}
\usepackage[T1]{fontenc}
\usepackage{graphicx}
\usepackage{longtable}
\usepackage{wrapfig}
\usepackage{rotating}
\usepackage[normalem]{ulem}
\usepackage{amsmath}
\usepackage{amssymb}
\usepackage{capt-of}
\usepackage{hyperref}
\usepackage{fancyhdr}
\usepackage{graphicx}
\usepackage{float}
\usepackage{booktabs}
\usepackage{indentfirst}
\usepackage[letterpaper, margin=2.5cm, top=3.5cm, bottom=2.5cm, headheight=2cm, footskip=1.2cm]{geometry}
\pagestyle{fancy}
\fancyhf{}
\lhead{\includegraphics[height=1.75cm]{../../../Images/logoepn2.png}}
\rhead{\fontsize{12}{14}\selectfont Escuela Politécnica Nacional \\ Facultad de Ingeniería de Sistemas \\ Compiladores y Lenguajes (ICCD422)\vspace{0.005cm}}
\fancyhfoffset{0.5cm}
\renewcommand{\headrulewidth}{0.4pt}
\usepackage{lastpage}
\fancyfoot{}
\fancyfoot[R]{\footnotesize Página \textbf{\thepage} de \textbf{\pageref{LastPage}}}
\renewcommand{\footrulewidth}{0.4pt}
\usepackage[T1]{fontenc}
\usepackage[utf8]{inputenc}
\usepackage[spanish, american]{babel}
\usepackage{csquotes}
\author{Jordy Cartagena}
\date{\today}
\title{}
\hypersetup{
 pdfauthor={Jordy Cartagena},
 pdftitle={},
 pdfkeywords={},
 pdfsubject={},
 pdfcreator={Emacs 29.3 (Org mode 9.6.15)}, 
 pdflang={English}}
\usepackage[style=ieee]{biblatex}
\addbibresource{/home/jordy/CollegeRep/Homework/Redes/Bibliografia/ref.bib}
\begin{document}

\begin{titlepage}
\centering
% Logo de la Institución3
\includegraphics[width=0.3\textwidth]{../../../Images/logoepn.jpg}\\[1cm]

% Nombre de la Institución
{\huge \textbf{Escuela Politécnica Nacional}} \\[0.5cm]

% Nombre de la Facultad/Departamento
{\Large \textbf{Facultad de Ingeniería de Sistemas}} \\[2cm]

% Título del Trabajo
{\Large \textbf{Informes LEX}} \\[0.25cm]
{\large Ejercicio #11} \\[1.5cm]

% Integrantes
{\Large Jordy Cartagena, Matias Monge, Esteban Valenzuela} \\[1cm]

% Información del Profesor
{\Large Ing. Rodrigo Fabian Chancusig MSc.} \\[2cm]

% Carrera y Materia
{\Large \textbf{Ciencia de Datos e Inteligencia Artificial}} \\[0.5cm]
{\Large \textbf{Compiladores y Lenguajes (ICCD422)}}\\[2cm]

% Fecha y Periodo
{\large GR2CD} \\
{\large 23 de Noviembre del 2026} \\
{\large Quito, Ecuador}
\end{titlepage}
\newpage
\setcounter{page}{1}

\section{Objetivo}
\label{sec:org2cd83c3}
\begin{itemize}
\item Implementar un analizador léxico que reconozca números binarios
(secuencia de 0 y 1) y muestre su valor decimal.
\end{itemize}

\section{Procedimiento y Ejecución}
\label{sec:orge72fc43}

La práctica constó en la implementación y ejecución del código de análisis léxico \texttt{lex}
(Ver el Anexo  \ref{anexo:a1}) que permita reconocer números
binarios, y puedan mostrar su valor decimal.

\appendix

\section{Anexo 1: Script \texttt{once.l}}
\label{sec:orga0551af}
\label{anexo:a1}
\end{document}
